% ============================================
% Pairwise Invasion Experiments
% ============================================
\subsection*{G. Pairwise Invasion Experiments}
\label{sec:suppl:invasion}

A brief summary of pairwise invasion experiments is provided in the main text Methods section. Here we provide extended experimental details.

% --------------------------------------------
% G.1 Experimental Design
% --------------------------------------------
\subsubsection*{G.1 Experimental Design}

To directly measure pairwise interactions between species, we performed reciprocal ``mutual invasion'' assays using the 12 most abundant isolates (ranked by ASV counts across all communities). Each pair of isolates was tested in both directions: species A as resident with species B as invader, and vice versa. This reciprocal design allows us to assess both competitive and facilitative interactions, as well as detect potential bistability where the outcome depends on initial conditions.

% --------------------------------------------
% G.2 Invasion Protocol
% --------------------------------------------
\subsubsection*{G.2 Invasion Protocol}

Each invasion assay was initiated by mixing isolate pairs at a 95:5 ratio (resident:invader) by colony count from equal-biomass cultures. The co-cultures were then subjected to 4 cycles of daily dilution ($\times$30 every 24 h) under the same growth conditions as community experiments (800 rpm shaking at 25$^\circ$C in 96-well deep-well plates). Final compositions were determined by colony counting on agar plates after the fourth dilution cycle.

Invasion outcomes were classified as:
\begin{itemize}
    \item \textbf{Coexistence}: Both species maintained $>$10\% relative abundance
    \item \textbf{Exclusion}: One species fell below 1\% relative abundance
    \item \textbf{Bistability}: Outcome depended on which species started as resident
\end{itemize}

% --------------------------------------------
% G.3 Nutrient Condition Effects
% --------------------------------------------
\subsubsection*{G.3 Nutrient Condition Effects}

Invasion assays were performed across all three nutrient conditions:
\begin{itemize}
    \item \textbf{Nutr$-$}: No added glucose or urea
    \item \textbf{Base}: 5 g/L glucose + 4 g/L urea
    \item \textbf{Nutr$+$}: 20 g/L glucose + 16 g/L urea
\end{itemize}

The frequency of competitive exclusion increased with nutrient availability, consistent with the coalescence outcome patterns observed in community experiments. Under Nutr$-$, more pairs achieved stable coexistence, while Nutr$+$ conditions led to more frequent exclusion of one species by the other.

